% Vietnamese translation for OI Public Library
% Source-File: IOI中国国家候选队论文/2004/林涛.ppt
\documentclass[11pt]{article}
\usepackage{oipl-vn}

\newcommand{\Oh}{\mathcal{O}}

\title{Ứng dụng của cây đoạn\\{\large Ghi chú theo slide}}
\author{Lin Tao}
\date{IOI 2004}

\begin{document}
\maketitle

\origtitle{Applications of segment trees}
\sourcepath{IOI 2004 / Lin Tao.ppt}

\section{Cấu trúc cơ bản}

Slide định nghĩa nút \(T(a,b)\) quản lý đoạn \([a,b]\). Khi độ dài lớn hơn
\(1\), nút tách thành hai nửa; khi độ dài bằng \(1\), ta tới lá. Vì chiều cao
cây là \(\Oh(\log L)\), các thao tác cập nhật và truy vấn trên đoạn thường
được đưa về một số nút chuẩn có kích thước \(\Oh(\log L)\).

\section{Các ví dụ}

Bài \emph{Snake} của SGU 128 xuất hiện như ví dụ đầu tiên: truy vấn hình chữ
nhật \(x_1<x<x_2\), \(y_1<y<y_2\) được xử lý bằng cách lưu thông tin theo
trục tọa độ.

Bài \emph{Cuboid} của POI 1999 có \(N\le 5000\) điểm \(P(x,y,z)\) và quan hệ
trội \(P_x>Q_x\), \(P_y>Q_y\), \(P_z>Q_z\). Các ký hiệu \(y[i]\),
\(\mathit{max}[i]\) và \(T[x,y]\) trên slide cho thấy cây đoạn được dùng để
duy trì giá trị tốt nhất theo một hoặc hai chiều sau khi đã sắp xếp chiều còn
lại.

Ví dụ thứ ba là hệ thống thống kê chiến trường: truy vấn và cập nhật trong
lưới \([1,2^k]\times[1,2^k]\). Cây đoạn hai chiều hoặc cây đoạn kết hợp mảng
có thể đổi bộ nhớ lấy thời gian, tùy giới hạn \(n,m,v\).

\section{Bài học}

Cây đoạn mạnh vì biến truy vấn liên tục trên đoạn thành hợp của ít nút rời
nhau. Khi bài toán có nhiều chiều, cần cân nhắc chiều nào sắp xếp ngoại tuyến,
chiều nào đưa vào cây, và thông tin nào thật sự phải lưu ở mỗi nút.

\end{document}
